\PassOptionsToPackage{table,dvipsnames}{xcolor}

\documentclass[sigconf,balance=false]{acmart}

\usepackage{popets}

\input{top-matter/header.tex}
\input{top-matter/macros.tex}

\newif\ifsubmission
\submissionfalse

% Copyright
\setcopyright{popets}
\copyrightyear{YYYY}

% Issue info
\acmYear{YYYY}
\acmVolume{YYYY}
\acmNumber{X}
\acmDOI{XXXXXXX.XXXXXXX}
\acmISBN{}
\acmConference{Proceedings on Privacy Enhancing Technologies}
\settopmatter{printacmref=false,printccs=false,printfolios=true}

\settopmatter{authorsperrow=4}

\begin{document}
\title{Beyond the Output: Inference Attacks on Private Set Union and Multi-Key Private Matching}

%%%%%%%%%%%%%%%% Authors' Info %%%%%%%%%%%%%%%%%
\author{Andrea Raguso}
%\orcid{1234-5678-9012}
\affiliation{%
  \institution{ETH Zurich}
  \city{Zurich}
  \state{}
  \country{Switzerland}}

\author{Francesca Falzon}
\affiliation{%
  \institution{ETH Zurich}
  \city{Zurich}
  \state{}
  \country{Switzerland}}

\author{Tianxin Tang}
\affiliation{%
  \institution{University of Glasgow}
  \city{Glasgow}
  \state{}
  \country{UK}}

\author{Kenneth G. Paterson}
\affiliation{%
  \institution{ETH Zurich}
  \city{Zurich}
  \state{}
  \country{Switzerland}}

\renewcommand{\shortauthors}{Raguso et al.}

\begin{abstract}
Recent work (Falzon and Tang USENIX 2025) has shown that a protocol participant who behaves honestly but strategically chooses its inputs can break input privacy in the Private Join and Compute functionality.
In this work, we expand our understanding of attacks in this setting by investigating a broader class of functionalities, namely: Private Set Union (PSU), PSU-Cardinality (PSU-CA), and Meta's multi-key private matching (MKPM) functionality.

We first present an attack on PSU that reconstructs the intersection using only two protocol invocations. We then show that any attack on PSI-Cardinality, such as that of Guo et al.\ (USENIX~2023), lifts to an attack on PSU-CA that recovers the intersection using only three additional queries. For the MKPM protocol, we distinguish its intended matching functionality from the additional protocol-specific leakage. We give an attack against the intended functionality, and then show that exploiting the additional leakage enables even stronger attacks, including partially reconstructing the other party’s records from a single protocol invocation.

We conclude by discussing possible mitigations for deploying such systems. Our analysis demonstrates limitations of existing secure multi-party computation security definitions and highlights the real-world privacy risks associated with deploying these functionalities in practice.
\end{abstract}


\keywords{private set union, private set intersection, multi-key private
  matching, attacks, multi-party computation}

\maketitle

\input{content/introduction.tex}
\input{content/prelims.tex}
\input{content/recovery-goals.tex}
\input{content/psu-attacks.tex}
\input{content/MRR.tex}
\input{content/VIR.tex}
\input{content/MSR.tex}
\input{content/evaluation.tex}
\input{content/conclusion.tex}

\section*{Ethical Considerations}
This work assumes intended use of the PSU, PSU-CA, and Meta’s MKPM protocols. Our analysis leverages intrinsic vulnerabilities in the functionality outputs and does not break the protocols themselves or the underlying cryptographic primitives. 

\heading{Use of Synthetic Data.}
Our experiments use only synthetic data, generated either by (a) sampling numerical values from integer intervals or (b) using the Python Faker library~\cite{faker}. This ensures that no personally identifying or sensitive data is leaked and eliminates any risk of harm to individuals or groups.

\heading{Disclosure.}
We contacted the authors of the MKPM work~\cite{MKPMC} on 27.02.2026, 
to share our findings and a preprint of this work. We received an official response on 02.03.2026 from the developers acknowledging the leakage of the protocol. They further stated that MKPM is not currently being used in any production systems.

Our analysis was thus conducted on the functionality of a protocol that, at the time of writing, is not deployed and therefore does not pose any negative impacts on live systems or stakeholders.

\section*{Open Science}
In accordance with the open science policy, our code base has undergone artifact evaluation. Our artifact contains the code for our attacks and for generating the synthetic data, as well as instructions on how to run the code. The open source respository can be found here: \href{https://github.com/deRaguso/mkpid-attacks-artifact}{https://github.com/deRaguso/mkpid-attacks-artifact}.


\begin{acks}
Francesca Falzon was supported by the Zurich Information Security and Privacy Center (ZISC).
Tianxin Tang was partially supported by an NWO VIDI grant (Project No. VI.Vidi.193.066).
The authors would also like to thank the reviewers for their helpful feedback, including their suggestion for simplifying the PSU attack.
The authors acknowledge the use of generative AI tools solely for basic editing
of the text and grammar. 
\end{acks}

\bibliographystyle{ACM-Reference-Format}
\bibliography{cryptobib/abbrev3,cryptobib/crypto,references}

\appendix
\input{content/ap-prelims.tex}
\subsection{Correctness conditions of $T$-Reconstruction}
\label{sec:correctness-TR}
Injectivity of $\varphi$ ensures that each $\by^* \in R_Y$ corresponds to
exactly one $\by \in Y$. Condition~(1) enforces per-record soundness, i.e.,
every value in a reconstructed record appears in the corresponding original
record. Condition~(2) captures $\Ispace[T]$-restricted per-record completeness
i.e., each reconstructed record contains all values from its original record
that also appear in $\Ispace[T]$. Finally, Condition~(3) enforces
$\Ispace[T]$-restricted global completeness i.e., all records containing a value
from $\Ispace[T]$ are at least partially reconstructed.

\input{content/ap-guos.tex}
\input{content/ap-psu.tex}
\input{content/ap-skpm.tex}
\section{Additional Content for MRR Attack}\label{ap:MRR}

\subsection{Proof of Theorem~\ref{thm:psu-mkpm:correctness}}
\label{sec:psu-mkpm:correctness}

In order to prove correctness, we must first prove the following two lemmas.:

\begin{lemma}
	\label{lem:set_isolation_determinism}
	Let $X, Y \subseteq \Ispace{}^{\leq \lambda}$ be two sets of records. 
	If $X$ is $Y$-isolated, 
	the cardinality of $\UID$ output by $\calF_\MKPM(X; Y)$ (Figure~\ref{fig:MKPM}) 
  is consistent across multiple evaluations.
\end{lemma}

\begin{proof}
	For any $\bx \in X$, let $Y_\bx\subseteq Y$ denote the set of records in $Y$
  that share at least one identifier with $\bx$,
	i.e., 
  \[
    Y_\bx := \{\by \in Y \setdsc \exists i, j \in \NN \text{ s.t. } \bx[i] = \by[j]\}.
  \]
	If $X$ is $Y$-isolated, then by definition, we have that $Y_\bx \cap Y_{\bx'} = \emptyset$ 
  for any $\bx, \bx' \in X$ such that $\bx \neq \bx'$. 
	Therefore, there exist no two records in $X$ that could be 
  matched with the same record in $Y$.
	
	We first consider sets $Y_\bx$ that are not empty.
	Since $\bx$ can only be matched with some $y \in Y_\bx$ and all $y \in Y_\bx$ 
  can only be matched with $\bx$,
	$\match$ (Figure~\ref{fig:match_logic}) will match $\bx$ with exactly one 
  $\by_\bx \in Y_\bx$, which are assigned the same UID 
  (line~\ref{lin:match_common_uid}).
	All other $\by \in Y_\bx\setminus\{\by_\bx\}$ will remain unmatched 
  and are assigned a unique UID (line~\ref{lin:match_unmatched_P_uid}).
  Furthermore, all $\by \in Y$ which do not belong to any $Y_{\bx'}$ 
  for some $\bx' \in X$ will also remain unmatched and are 
  assigned a unique UID (also line~\ref{lin:match_unmatched_P_uid}).
	Lastly, all $\bx\in X$ for which $Y_\bx = \emptyset$ will remain 
  unmatched and receive a distinct UID (line~\ref{lin:match_unmatched_C_uid}). 
	
	Thus, $\UID$ contains one UID for each $\by\in Y$ and one UID for 
  each unmatched $\bx\in X$, i.e.,
	$|\UID| = |Y| + |\{\bx \in X \setdsc Y_\bx = \emptyset\}|$.
	This is independent of the order of the records in $X$ and $Y$ and
	also independent of the shuffling done before the matching step in $\MKPM$.
	Since $\MKPM$ only uses randomness to shuffle the inputs, the lemma follows.
\end{proof}

\begin{lemma}
  \label{lem:set_isolation_inference}
  Let $X$ and $V\subseteq \Ispace{}^{\leq \lambda}$ be two sets of records and
  let $X_1$ and $X_2$ be a partition of $X$.  Moreover, let
  $(\UID, \MX) \sample \calF_{\MKPM(X; Y)}$ and
  $(\UID_i, M_{X,i}) \sample \calF_{\MKPM(X_i; Y)}$ for $i \in \{1,2\}$.  If $X$
  is $V$-isolated, we have $|\UID| = |\UID_1| + |\UID_2| - |V|$.
\end{lemma}

% \begin{proof}
% 	We show this by induction on $|\setX_1|$.
% 	If $|\setX_1| = 0$ all records of $\setY$ are unmatched and, thus, $|\UID_1| = |Y|$.
% 	Note that $\setX_2 = \setX$ and thus we have $|\UID| = |\UID_2| = |\UID_2| + |\UID_1| - |Y|$.
% 	In the first equality, we use \cref{lem:set_isolation_determinism}. 

% 	Assume $|\setX_1| = n_1$ for some $n_1 > 0$ and $|\UID| = |\UID_1| + |\UID_2| - |\setY|$.
% 	Since $\setX_1$ and $\setX_2$ form a partition of $\setX$, we must move a record from $\setX_2$ to $\setX_1$
% 	to achieve $|\setX_1| = n_1+1$. Note that $|\UID|$ remains unchanged due to \cref{lem:set_isolation_determinism}.
% 	Let $x \in \setX_2$ and let $\setX'_1 := \setX_1 \cup \{x\}$ and $\setX'_2 := \setX_2 \setminus \{x\}$.
% 	Moreover, let $\UID'_1$ and $\UID'_2$ be the sets of UIDs resulting from evaluating $\MKPM(\setX'_1, \setY)$ and $\MKPM(\setX'_2, \setY)$. 
% 	We distinguish two cases.
% 	If no identifier of $x$ occurs in $\setY$, $x$ was assigned its own UID.
% 	Therefore, $|\UID'_1| = |\UID_1| + 1$ and $|\UID'_2| = |\UID_2| - 1$, which implies the claim.
	
% 	For the second case, assume $x$ shares some identifiers with $n^*$ records of $\setY$.
% 	Call this set $\setY_x$.
% 	Since \setX{} is \setY-isolated, all records $y \in \setY_x$ only share identifiers with $x$,
% 	but no other $x'\in \setX_2$.
% 	Thus, $x$ is matched with some $y^* \in \setY_x$,
% 	i.e., $x$ and $y^*$ are assigned the same $\uid \in \UID_2$.
% 	After removing $x$ from $\setX_2$, $y^*$ will still be assigned some $\uid' \in \UID'_2$, 
% 	which it does not share with any $x' \in \setX'_2$, again since $\setX$ is \setY-isolated. 
% 	Therefore, $|\UID_2| = |\UID'_2|$. 
% 	Similarly, since $x\not\in \setX_1$, no $y\in \setY_x$ and $x' \in \setX_1$ are assigned the same $\uid \in \UID_1$.
% 	By the definition of $\setY_x$ and since \setX{} is \setY-isolated, $x$ will be assigned the same $\uid\in \UID'_1$
% 	as some $y\in\setY_x$ after being added to $\setX_1$, implying $|\UID_1| = |\UID'_1|$.
% 	The claim then follows trivially.
% \end{proof}


We can now use the above lemmas to finally prove Theorem~\ref{thm:psu-mkpm:correctness}.


% \begin{proof}
% 	Let \posSet{} and \negSet{} be the two sets output by $\PSUCAattack^{\MKPM(\cdot, \setV)}(\setT)$.
% 	Note all records of \setT{} are contained in some subset $\setT_c\subseteq \setT$ that will eventually reach line~\ref{lin:PSUCA_add_condition},
% 	since the attack only terminates once the priority queue is empty. 
% 	Therefore, every record of \setT{} is added to either \posSet{} or \negSet{}.
% 	For any $\setT_c$ reaching line~\ref{lin:PSUCA_add_condition} 
% 	we have that either (1) $k_c = |\setT_c|$ or (2) $k_c  = 0$ by the loop condition on line~\ref{lin:PSUCA_inner_while_condition}.
% 	Let $\UID_c$ result from evaluating $\MKPM(\setT_c, \setV)$.
% 	By \cref{eqn:union_intersect_cardinalities} and \cref{lem:set_isolation_inference}, we have $k_c = |\setT_c| + |\setV| - |\UID_c|$.
% 	\begin{description}
% 		\item[Case (1)] If $k_c = |\setT_c|$, then we have $|\UID_c| = |\setV|$, i.e., all records of $\setT_c$ are matched. 
% 		Since records are only added to \posSet{} if this case applies, we have $\posSet \subseteq \setT \sqcap \setV$.
% 		\item[Case (2)] If $k_c = 0$, then $|\UID_c| = |\setT_c| + |\setV|$, i.e., no records of $\setT_c$ 
% 		were matched. Therefore, no record in $\setT_c$ shares any identifiers with any record in $\setV{}$.
% 		Since records are only added to \negSet{} if this case applies, we have $\negSet \subseteq \setT \setminus (\setT \sqcap \setV)$. 
% 	\end{description}
% 	We have therefore shown that \posSet{} only contains matchable records, \negSet{} only contains non-matchable records 
% 	and that every record in \setT{} is contained in either \posSet{} or \negSet{}.
% 	This proves the theorem.
% \end{proof}
\input{content/ap-VIR.tex}
\section{Proofs for Analysis of $\calF_{\LMKPM}$}

\subsection{Proof of Correctness of $\dictInfer$}
\label{sec:dictinfer:correctness}


\begin{lemma}
	\label{lem:dictinfer_MR}
	Let $X$ and $Y$ be two sets of records and let $(\hXone, \hYone)$ be the
        hidden shuffled records received by $\tP_1$.  Moreover, let
        $f_1: \Ispace{} \to \Uspace$ be the hiding function sampled in
        $\calF_\LMKPM$ to produce $\hXone$ and $\hYone$.  If $D$ is a dictionary
        mapping $f_1(v)$ to $v$ for all $v \in \Ispace[X]$, then
        $\dictInfer(D, \hYone)$ (Figure~\ref{fig:dict_infer}) produces an
        $X$-reconstruction of $Y$.
\end{lemma}
\begin{proof}
	Let $R$ be the multiset produced by \dictInfer.
	Note that for each $\hby \in \hYone$ exactly one reconstructed record
        $\br$ is added to $R$. Let $d: R \to \hYone$ denote the function mapping
        these reconstructed records to the corresponding $\hby$. Clearly, $d$ is
        a bijection.

	We now define an injective function $\varphi: R \to Y$ and prove that it
        satisfies the three conditions for $X$-Reconstruction (Definition
        \ref{def:max_recons}). For any record $\br \in R$ let
        $(d_1, \dots, d_k)=d(\br) \in \hYone$ and let
	$$\varphi(\br) := (f_1^{-1}(d_1), \dots, f_1^{-1}(d_k))$$
	
	Note that $f_1^{-1}(d_j)$ exists for all $\br$ and $j$, since $d(\br) \in \hYone$ by definition.
	Moreover, since $f_1$ is injective, $f_1^{-1}$ is injective for elements in $f_1(\Ispace)$.
	$\varphi$ is injective since $d$ and $f_1^{-1}$ are injective.
	We show the three properties from Definition~\ref{def:max_recons} separately:
	\begin{enumerate}
		\item Let $\br \in R$, $k := |\br|$ and $i \leq |\br|$. 
		Moreover, let $(d_1, \dots, d_{k'}) = d(\br)$ for some $k' \geq k$.
		By definition of $d$ and by construction of $\br$ (line~\ref{lin:substitute_add_id}), 
		there exists some $j^*\in[k']$ such that $r[i] = D[d_{j^*}]$.
		Suppose for a contradiction that $\br[i] \not \in \varphi(\br)$, that is,
		we have $\br[i] \neq f_1^{-1}(d_j)$ for all $j \in [k']$.
		But since for all $j \in [k']$ with $d_j \in D$ we have $f_1^{-1}(d_j) = D[d_j]$,
		we also have that $\br[i] \neq D[d_j]$, i.e., there is no $j$ such that $\br[i] = D[d_j]$. 
		This is a direct contradiction to $\br[i] = D[d_{j^*}]$.
		\item Let $\br\in R$ and $(d_1, \dots, d_{k'}) = d(\br)$ for some $k' \in \N$.
		Let $v \in \varphi(\br)$, 
		i.e., there is some $i \in [k']$ such that $v = f_1^{-1}(d_i)$.
		If $v \in \Ispace[X]$, then by assumption we have $f_1(v) \in D$ and $D[f_1(v)] = v$.
		Since $v = f_1^{-1}(d_i)$, we have $f_1(v) = f_1(f_1^{-1}(d_i)) = d_i$ 
		and thus $d_i \in D$. Therefore, the condition on line~\ref{lin:substitute_membership_check} is met and
		$D[d_i] = D[f_1(v)] = v$ is added to $\br$.
		\item Holds trivially, since we have $\varphi(R) = Y$ by the definition of $d$. 
	\end{enumerate}
The lemma thus follows.
\end{proof}

\subsection{Pseudocode of $\recenumattack$}
\label{sec:recenumattack:pseudocode}
The pseudocode of $\recenumattack$ can be found in Figure~\ref{fig:rec_enum_attack}.


\subsection{Proof of Correctness of $\recenumattack$}
\label{sec:recenumattack:correctness}
\begin{proof}
In order to prove the theorem we prove the following proposition. 
The theorem follows directly from Proposition~\ref{thm:recenum_D} and Lemma~\ref{lem:dictinfer_MR}.
\begin{proposition}\label{thm:recenum_D}
	Let $T$ and $Y$ be two sets of records and let
	$D$ be the dictionary constructed during an execution of\newline$\recenumattack^{\calF_\LMKPM(\cdot, Y)}(T)$ (Figure~\ref{fig:rec_enum_attack}).
	Moreover, let $f_1$ be the hiding function sampled during the evaluation of $\calF_\LMKPM$ in line~\ref{lin:recenum_eval} in Figure~\ref{fig:rec_enum_attack}. 
	We have $D[f_1(v)] = v$ for all $v \in \Ispace[T]$.
\end{proposition}
\begin{proof} \textbf{of Proposition~\ref{thm:recenum_D}} 
	Let $X = \{\bx_0, \dots, \bx_{n-1}\}$, where $\bx_i$ denotes the altered target record constructed in 
	line~\ref{lin:recenum_extend} for $i \in \{0, \dots, n - 1\}$.
	We show that $D[f_1(v)] = v$ for all $v \in \Ispace[X]$.
	Since $\Ispace[T] \subseteq \Ispace[X]$, this implies the claim.
	
	We first show that $u_0 = f_1(v_0)$ for the value $u_0$ defined on line~\ref{lin:recenum_e0}.
	Note that the transformation of an index's binary representation to a sequence of $v_0$ and $v_1$
	in lines \ref{lin:idx_to_seq_bin_encode_start}~-~\ref{lin:idx_to_seq_bin_encode_end} of \IdxToSeq{}
	could produce two encodings that only contain one value, namely, $(v_0)^\ell$ and $(v_1)^\ell$
	for $i = 0$ and $i = 2^\ell - 1$ respectively.
	However, the latter is prevented by the condition on line~\ref{lin:seq_to_idx_tiebreak} and therefore
	\IdxToSeq{} only outputs one sequence where all values are equal.
	Hence, there is only one record in $X$ whose first $\ell$ values all coincide,
	namely, $\bx_0$ with encoding $(v_0)^\ell$.
	Since $f_1$ is injective and the values in records of $X$ are not shuffled, 
	the latter also holds for $\hXone$.
	Therefore, $\hbx_0 = (f_1(\bx_0[1]), \dots, f_1(\bx_0[\lambda + \ell]))$ (line~\ref{lin:recenum_find0}).
	Since $\bx_0[1] = v_0$, we have $u_0 = \hbx_0[1] = f_1(\bx_0[1]) = f_1(v_0)$ (line~\ref{lin:recenum_e0}).
	
	We now show that the attack correctly recovers indices from $\hXone$. 
	We have already established that $\hbx_0$ is correctly matched with $\bx_0$,
	as it's encoding is the only one containing only $v_0$. 
	Any other encoding $\bs$ produced by \IdxToSeq{}
	either (a) is $(v_1)^{\ell - 1} \| v_2$ or (b) only
	contains the values $v_0$ and $v_1$.
	\begin{description}
	\item[Case (a):] This case is only produced on input $i = 2^\ell - 1$ and the resulting vector of hidden values in the 	
		 leakage is $(f_1(v_1))^{\ell - 1} \| f_1(v_2)$. 
		 This satisfies the condition on line~\ref{lin:seq_to_idx_tiebreak} of \SeqToIdx{}, 
		 which therefore yields the correct index $2^\ell - 1$.
	\item[Case (b):] In this case, the vector of hidden values corresponding to $\bs$ in the leakage, which we denote
		$\hbs:=(f_1(\bs[1]),\allowbreak\dots,f_1(\bs[\ell]))$, does not satisfy the condition on line~\ref{lin:seq_to_idx_tiebreak}, 
		since the only encoding that would do so is $(v_1)^{\ell}$. 
		But this is never output by \IdxToSeq{}. 
		Therefore, $\hbs$ is interpreted as a binary encoding, where every occurrence of $u_0$ 
		is interpreted as $0$ and every other hidden value as $1$.
		The correctness of this recovery therefore follows from the fact that $u_0 = f_1(v_0)$,
		which we established above.
	\end{description}
	
	By the definition of $\calF_\LMKPM$, we have $\hbx = (f_1(\bx_i[1]), \dots, f_1(\bx_i[\lambda+\ell]))$ for all $i \in \{0, \dots, n-1\}$.
	Since $\Ispace[X] := \bigcup_{\substack{l}{\bx\in X\\v \in \bx}} v$, 
	the claim holds by the construction of $D$ in lines~\ref{lin:recenum_dict_start}~-~\ref{lin:recenum_dict_end}.
\end{proof}
	
% \begin{proof}
	The theorem follows directly from Proposition~\ref{thm:recenum_D} and Lemma~\ref{lem:dictinfer_MR}.
\end{proof}

\subsection{Pseudocode of $\snakeattack$}
\label{sec:snakeattack:pseudocode}

The pseudocode of $\recenumattack$ can be found in Figure~\ref{fig:snakeattack}.



\subsection{Proof of Correctness of $\snakeattack$}
\label{sec:snakeattack:correctness}


\begin{proof}
	To prove the theorem, we prove the following proposition.
	The theorem follows directly from Proposition~\ref{thm:snake_D} and Lemma~\ref{lem:dictinfer_MR}.
\begin{proposition}
	\label{thm:snake_D}
	Let $T$ and $Y$ be two sets of records and let $D$ be the substitution dictionary
	constructed during an invocation of $\snakeattack^{\calF_\LMKPM(\cdot, Y)}(T)$.
	Moreover, let $f_1$ be the hiding function sampled during the invocation of $\calF_\LMKPM$
	in line~\ref{lin:snake_evaluation}.
	Then, we have $D[f_1(v)] = v$ for all $v \in \Ispace[T]$.
\end{proposition}
\begin{proof} \textbf{of Proposition~\ref{thm:snake_D}}
	Let $X := \{\bx_1, \dots, \bx_n\}$ where $\bx_i$ is produced by \snakeencodeUnique{} in line~\ref{lin:snake_encode}
	and let $k' := |\bx_i|$ for any $i\in [n]$\footnote{We assume that all records have the same length $\lambda$.
	Thus, $k' \in \{\lambda + 1, \lambda+2\}$}.
	We show that $D[f_1(v)] = v$ for all $v \in \Ispace[X]$.
	Since \snakeattack{} potentially adds fresh values to records of $T$, but does not remove any, 
	we have $\Ispace[T]\subseteq\Ispace[X]$, which implies the claim.
	
	Similar to the proof of Proposition~\ref{thm:recenum_D}, we show that
	$\hbx_i[j] = f_1(\bx_i[j])$ for all $i\in[n]$ and all $j \in [k']$.
	The result then follows from the construction of $D$ in lines~\ref{lin:snake_D_start}~-~\ref{lin:snake_D_end}. 
	In particular, we will show that the linked list we create in \IdxToSeq{} is preserved in the leakage, 
	and that this linked list allows us to uniquely identify each record. 
	We show this by induction on $i$. 
	Let $j^*$ denote the index of a column whose values are all unique (note that if such a column does not exist, it is created on line~\ref{lin:distinctIDs_j*}).

	For $i = 1$, observe that the value $t_1[j^*]$ is not appended to any record in \snakeencodeUnique{}; 
	in other words, it is the only item in the linked list that doesn't have anything pointing to it.
	In contrast, for all $i' > 1$, the value $t_{i'}[j^*]$ is appended to $t_{i'+1}$.
	$\hbx_1$ is therefore the only record in $X$ whose $j^*$-th value is not contained in the last column, 
	i.e., $\hbx_1[j^*] \not \in \{\hbx_{i'}[k'] \setdsc i'\in[n]\}$.
	By the definition of $\calF_\LMKPM$, this property is preserved in the leakage $\hXone$.
	That is, the corresponding record in $h = (f_1(\hbx_1[1]), \dots, f_1(\hbx_1[k']))$
	is the only record in $\hXone$ such that $h[j^*]\not\in\{h_{i'}[k'] \setdsc i' \in [n]\}$.
	In line~\ref{lin:snake_recover_c} in \snakerecoverUnique{} we pick $h_\ind := h$.
	Thus, for all $j \in [k']$ we have that $h_\ind[j] = f_1(\bx_1[j])$
	and, therefore  $\hbx_1[j] = h_\ind[j] = f_C(\bx_1[j])$, where the first equality follows from line~\ref{lin:snake_recover_assign_t_star}.

	Now assume that for some $i < n$ we have $\hbx_i[j] = f_1(\bx_i[j])$ for all $j \in [k']$. 
	We show that there is a unique record which $\bx_i[k']$ is pointing to (that has not been traversed by the previous items in the linked list), 
	and that, once again this property is preserved in the leakage allowing for unique identification of the record.
	By construction in \snakeencodeUnique{}, we have $\bx_{i+1}[j^*] = \bx_i[k']$.
	Since all values in the $j^*$-th column are distinct, 
	$\bx_{i+1}$ is the only record in $X$ that satisfies this condition.
	As before, since $f_1$ is injective, the corresponding record in the leakage 
	$h = (f_1(\bx_{i+1}[1]), \dots, f_1(\bx_{i+1}[k']))$
	is the only record that satisfies $h[j^*] = \hbx_i[k']$.
	Moreover, all hidden values in the $j^*$-th column of $\hXone$	are distinct. % $t'_{i+1}$ is the only record in $\setT'$
	Thus, the dictionary $\loc$ in \SeqToIdx{} contains exactly the $n$ keys $\{f_1(\bx_{i'}[j^*]) \setdsc i' \in [n]\} = \{f_1(\hbx_{i'}[j^*]) \setdsc i' \in [n]\}$,
	which are mapped as $\loc[f_1(\bx_{i'}[j^*])] = i'$.
	In line~\ref{lin:snake_recover_assign_t_star}, we pick $\hbx_{i+1} := h_\ind$ for $\ind = \loc[\hbx_i[k']]$, 
	implying $\hbx_{i+1}[j^*] = \hbx_i[k']$ and, thus, $\hbx_{i+1} = h$.
	We therefore have $\hbx_{i+1}[j] = h[j] = f_1(\bx_{i+1}[j])$ for all $j \in [k']$.
\end{proof}

	The statement follows from Proposition~\ref{thm:snake_D} and Lemma~\ref{lem:dictinfer_MR}.
\end{proof}

\input{content/ap-evaluation.tex}
\end{document}